% GENERATED FILE. Edit canonical lessons, then run npm run generate:books.
% canonical-source-hash: fnv1a64:4967618c67c32abd
% canonical-lessons: TE-C369-rajabhavanam, TE-C369-ratham, TE-C369-varasatvam, TE-C369-kulcu, TE-C369-nirmincu

\chapter{Palaces and Chariots}
\label{ch:te-palaces-and-chariots}
\hlchaptermodality{369}

\begin{chapteropening}
\textbf{By the end of this chapter:} \emph{I can say a palace, a chariot, heritage, to demolish and to build.}

Complete the last lesson of chapter 369: I can say a palace, a chariot, heritage, to demolish and to build.
\end{chapteropening}

\section[\texorpdfstring{\te{రాజభవనం} (rājabhavanaṁ)}{రాజభవనం (rājabhavanaṁ)}]{\te{రాజభవనం} (rājabhavanaṁ) — a palace}

\label{lesson:TE-C369-rajabhavanam}



\begin{quote}
Before the new one: say the Telugu for a civilisation, then the Telugu for an era.
\end{quote}

\subsection*{You'll want to know: \te{రాజభవనం}}
\textbf{\te{రాజభవనం}} — \emph{rājabhavanaṁ} — \textquotedblleft{}a palace\textquotedblright{}.

\begin{quote}
\te{రాజభవనంలో} \te{వంద} \te{గదులు} \te{ఉండేవి}.
\end{quote}

\begin{quote}
\emph{rājabhavanaṁlō vanda gadulu uṇḍēvi.}
\end{quote}

\begin{quote}
The palace had a hundred rooms.
\end{quote}

\te{రాజ}, king, with \te{భవనం}, a building.

\begin{cousinweb}
\textbf{\te{రాజభవనం}} (\emph{rājabhavanaṁ}) is borrowed from Sanskrit: \emph{rājan}, \textquotedblleft{}a king\textquotedblright{} (CDIAL 10679), and \emph{bhavana}, \textquotedblleft{}a place of abode\textquotedblright{} (CDIAL 9417).
\end{cousinweb}

\begin{grammarlens}[title={What you've built}]
One more word for telling what happened.
\end{grammarlens}

\subsection*{Your turn}
\begin{itemize}
\raggedright
  \item \emph{Say it:} \emph{rājabhavanaṁ}
  \item \emph{Say it:} \emph{rājabhavanaṁ}, once more
  \item \emph{Say it:} \emph{rājabhavanaṁlō vanda gadulu uṇḍēvi.} — The palace had a hundred rooms.
  \item \emph{Recall:} say the Telugu for a civilisation, then the Telugu for an era, then say \emph{rājabhavanaṁ} again
\end{itemize}

\begin{tcolorbox}[breakable,colback=teal!4,colframe=teal!35!black,title={Before you move on}]
What does \te{రాజభవనం} mean? (\textquotedblleft{}A palace\textquotedblright{}.) Say it once more.
\end{tcolorbox}
\section[\texorpdfstring{\te{రథం} (rathaṁ)}{రథం (rathaṁ)}]{\te{రథం} (rathaṁ) — a chariot}

\label{lesson:TE-C369-ratham}



\begin{quote}
Before the new one: say the Telugu for an era, then the Telugu for a palace.
\end{quote}

\subsection*{You'll want to know: \te{రథం}}
\textbf{\te{రథం}} — \emph{rathaṁ} — \textquotedblleft{}a chariot\textquotedblright{}.

\begin{quote}
\te{రాజు} \te{రథం} \te{మీద} \te{యుద్ధానికి} \te{వెళ్ళాడు}.
\end{quote}

\begin{quote}
\emph{rāju rathaṁ mīda yuddhāniki veḷḷāḍu.}
\end{quote}

\begin{quote}
The king went to war in a chariot.
\end{quote}

Plural: \te{రథాలు}, \emph{rathālu}.

\begin{cousinweb}
\textbf{\te{రథం}} (\emph{rathaṁ}) is borrowed from Sanskrit \emph{ratha}, \textquotedblleft{}a chariot, a cart\textquotedblright{} (CDIAL 10602).
\end{cousinweb}

\begin{grammarlens}[title={What you've built}]
One more word for telling what happened.
\end{grammarlens}

\subsection*{Your turn}
\begin{itemize}
\raggedright
  \item \emph{Say it:} \emph{rathaṁ}
  \item \emph{Say it:} \emph{rathaṁ}, once more
  \item \emph{Say it:} \emph{rāju rathaṁ mīda yuddhāniki veḷḷāḍu.} — The king went to war in a chariot.
  \item \emph{Recall:} say the Telugu for an era, then the Telugu for a palace, then say \emph{rathaṁ} again
\end{itemize}

\begin{tcolorbox}[breakable,colback=teal!4,colframe=teal!35!black,title={Before you move on}]
What does \te{రథం} mean? (\textquotedblleft{}A chariot\textquotedblright{}.) Say it once more.
\end{tcolorbox}
\section[\texorpdfstring{\te{వారసత్వం} (vārasatvaṁ)}{వారసత్వం (vārasatvaṁ)}]{\te{వారసత్వం} (vārasatvaṁ) — heritage; inheritance}

\label{lesson:TE-C369-varasatvam}



\begin{quote}
Before the new one: say the Telugu for a palace, then the Telugu for a chariot.
\end{quote}

\subsection*{You'll want to know: \te{వారసత్వం}}
\textbf{\te{వారసత్వం}} — \emph{vārasatvaṁ} — \textquotedblleft{}heritage; inheritance\textquotedblright{}.

\begin{quote}
\te{ఆ} \te{పురాతన} \te{కట్టడాలు} \te{మన} \te{వారసత్వం}.
\end{quote}

\begin{quote}
\emph{ā purātana kaṭṭaḍālu mana vārasatvaṁ.}
\end{quote}

\begin{quote}
Those ancient buildings are our heritage.
\end{quote}

\te{వారసత్వంగా}, by inheritance: \te{వారసత్వంగా} \te{వచ్చిన} \te{ఇల్లు}, an inherited house.

\begin{cousinweb}
\textbf{\te{వారసత్వం}} (\emph{vārasatvaṁ}) is built with the Sanskrit ending \emph{-tva}, \textquotedblleft{}-ness\textquotedblright{}, on \emph{vārasu}, \textquotedblleft{}an heir\textquotedblright{}; the origin of \emph{vārasu} is not settled.
\end{cousinweb}

\begin{grammarlens}[title={What you've built}]
One more word for telling what happened.
\end{grammarlens}

\subsection*{Your turn}
\begin{itemize}
\raggedright
  \item \emph{Say it:} \emph{vārasatvaṁ}
  \item \emph{Say it:} \emph{vārasatvaṁ}, once more
  \item \emph{Say it:} \emph{ā purātana kaṭṭaḍālu mana vārasatvaṁ.} — Those ancient buildings are our heritage.
  \item \emph{Recall:} say the Telugu for a palace, then the Telugu for a chariot, then say \emph{vārasatvaṁ} again
\end{itemize}

\begin{tcolorbox}[breakable,colback=teal!4,colframe=teal!35!black,title={Before you move on}]
What does \te{వారసత్వం} mean? (\textquotedblleft{}Heritage; inheritance\textquotedblright{}.) Say it once more.
\end{tcolorbox}
\section[\texorpdfstring{\te{కూల్చు} (kūlcu)}{కూల్చు (kūlcu)}]{\te{కూల్చు} (kūlcu) — to demolish; to bring down}

\label{lesson:TE-C369-kulcu}



\begin{quote}
Before the new one: say the Telugu for a chariot, then the Telugu for heritage.
\end{quote}

\subsection*{You'll want to know: \te{కూల్చు}}
\textbf{\te{కూల్చు}} — \emph{kūlcu} — \textquotedblleft{}to demolish; to bring down\textquotedblright{}.

\begin{quote}
\te{వాళ్ళు} \te{పాత} \te{భవనాన్ని} \te{కూల్చారు}.
\end{quote}

\begin{quote}
\emph{vāḷḷu pāta bhavanānni kūlcāru.}
\end{quote}

\begin{quote}
They demolished the old building.
\end{quote}

Past (I …): \te{కూల్చాను}, \emph{kūlcānu}.

The causative of \te{కూలు}, to fall: to make something fall.

\begin{cousinweb}
\textbf{\te{కూల్చు}} (\emph{kūlcu}) is native Dravidian: DEDR 1907 lists Telugu \emph{kūlcu}, \textquotedblleft{}to cause to fall, knock down\textquotedblright{}, beside \emph{kūlu}, \textquotedblleft{}to fall down\textquotedblright{}.
\end{cousinweb}

\begin{grammarlens}[title={What you've built}]
One more word for telling what happened.
\end{grammarlens}

\subsection*{Your turn}
\begin{itemize}
\raggedright
  \item \emph{Say it:} \emph{kūlcu}
  \item \emph{Say it:} \emph{kūlcu}, once more
  \item \emph{Say it:} \emph{vāḷḷu pāta bhavanānni kūlcāru.} — They demolished the old building.
  \item \emph{Recall:} say the Telugu for a chariot, then the Telugu for heritage, then say \emph{kūlcu} again
\end{itemize}

\begin{tcolorbox}[breakable,colback=teal!4,colframe=teal!35!black,title={Before you move on}]
What does \te{కూల్చు} mean? (\textquotedblleft{}To demolish; to bring down\textquotedblright{}.) Say it once more.
\end{tcolorbox}
\section[\texorpdfstring{\te{నిర్మించు} (nirmiñcu)}{నిర్మించు (nirmiñcu)}]{\te{నిర్మించు} (nirmiñcu) — to build; to construct}

\label{lesson:TE-C369-nirmincu}



\begin{quote}
Before the new one: say the Telugu for heritage, then the Telugu for to demolish.
\end{quote}

\subsection*{You'll want to know: \te{నిర్మించు}}
\textbf{\te{నిర్మించు}} — \emph{nirmiñcu} — \textquotedblleft{}to build; to construct\textquotedblright{}.

\begin{quote}
\te{ఆ} \te{రాజు} \te{ఒక} \te{పెద్ద} \te{కోట} \te{నిర్మించాడు}.
\end{quote}

\begin{quote}
\emph{ā rāju oka pedda kōṭa nirmiñcāḍu.}
\end{quote}

\begin{quote}
That king built a great fort.
\end{quote}

Past (I …): \te{నిర్మించాను}, \emph{nirmiñcānu}.

A formal partner of \te{కట్టు}; the noun is \te{నిర్మాణం}, construction.

\begin{cousinweb}
\textbf{\te{నిర్మించు}} (\emph{nirmiñcu}) is built on Sanskrit \emph{nirmāṇa}, \textquotedblleft{}measuring; making\textquotedblright{} (CDIAL 7368), with the Telugu verb ending \emph{-iñcu}.
\end{cousinweb}

\begin{grammarlens}[title={What you've built}]
One more word for telling what happened.
\end{grammarlens}

\subsection*{Your turn}
\begin{itemize}
\raggedright
  \item \emph{Say it:} \emph{nirmiñcu}
  \item \emph{Say it:} \emph{nirmiñcu}, once more
  \item \emph{Say it:} \emph{ā rāju oka pedda kōṭa nirmiñcāḍu.} — That king built a great fort.
  \item \emph{Recall:} say the Telugu for heritage, then the Telugu for to demolish, then say \emph{nirmiñcu} again
\end{itemize}

\begin{tcolorbox}[breakable,colback=teal!4,colframe=teal!35!black,title={Before you move on}]
What does \te{నిర్మించు} mean? (\textquotedblleft{}To build; to construct\textquotedblright{}.) Say it once more.
\end{tcolorbox}
